%% assembly -- the main theorem, stated once; counterexample words; the
%% decision procedure as a numbered algorithm with its dependency table;
%% countermodels from counterexample words; counterexample bounds per case;
%% the proof of the main theorem by assembling the earlier sections; and plain
%% entailment as a corollary. The remark on complexity is in complexity.tex.

\providecommand{\parity}{\operatorname{par}}
\providecommand{\Mod}{\operatorname{Mod}}
\providecommand{\CondI}{\textup{(I0)}}
\providecommand{\CondU}{\textup{(U)}}
\providecommand{\CondII}{\textup{(II)}}
\providecommand{\DeadT}{\mathsf{DeadT}}
\providecommand{\DeadU}{\mathsf{DeadU}}
\providecommand{\Fup}[1]{\Fact^{\uparrow}_{#1}}
\providecommand{\Bsrc}{\mathcal B}
\providecommand{\RU}{\mathcal R_{\mathrm U}}
\providecommand{\Lnm}{\mathsf L_{\mathrm{nm}}}
\providecommand{\Eall}{E_{\mathrm{all}}}
\providecommand{\Ereg}{E_{\mathrm{reg}}}
\providecommand{\Efin}{E_{\mathrm{fin}}}
\crefformat{asstep}{Step~#2#1#3}
\Crefformat{asstep}{Step~#2#1#3}
\crefrangeformat{asstep}{Steps~#3#1#4 to~#5#2#6}
\Crefrangeformat{asstep}{Steps~#3#1#4 to~#5#2#6}
\crefmultiformat{asstep}{Steps~#2#1#3}{ and~#2#1#3}{, #2#1#3}{, and~#2#1#3}
\Crefmultiformat{asstep}{Steps~#2#1#3}{ and~#2#1#3}{, #2#1#3}{, and~#2#1#3}
\crefrangemultiformat{asstep}{Steps~#3#1#4 to~#5#2#6}{ and~#3#1#4 to~#5#2#6}{, #3#1#4 to~#5#2#6}{, and~#3#1#4 to~#5#2#6}
\Crefrangemultiformat{asstep}{Steps~#3#1#4 to~#5#2#6}{ and~#3#1#4 to~#5#2#6}{, #3#1#4 to~#5#2#6}{, and~#3#1#4 to~#5#2#6}

\section{The Main Theorem}\label{sec:as}

This section states the main theorem and proves it by assembling the results
of \cref{sec:bg,sec:lf,sec:ch,sec:tc,sec:sm,sec:un,sec:ex,sec:nm,sec:mo}.
Apart from bookkeeping, only two things are proved here: how a
counterexample word yields a countermodel (\cref{lem:as-countermodel}), and
two short estimates that turn the length bounds of the earlier sections into
closed forms (\cref{prop:as-bounds}).

Throughout, the arity $n\ge0$ and the variance $c$ are fixed. For a flat
system $\varphi$ over $X$ and a flat system $\psi$ over $X\cup Z$ we write
\[
 v=\len X,\qquad h=\len Z,\qquad s=v+h,
\]
so that $s$ is the size parameter of \cref{sec:bg-constraints}, counted after
flattening. As in \cref{sec:mo-short}, $E_1(x)=2^x$ and
$E_{k+1}(x)=2^{E_k(x)}$.

\subsection{Statement}\label{sec:as-statement}

When the entailment fails although both systems are satisfiable, the failure
is certified by a single word, which comes from one of the three conditions
of \cref{thm:ch-characterization}.

\begin{definition}[Counterexample words]\label{def:as-certificate}
Let $\varphi$ over $X$ and $\psi$ over $X\cup Z$ be flat systems over $\BFT$,
both satisfiable. A \emph{counterexample word} of $(\varphi,\psi)$ is a word of
one of the following three kinds, together with the signed variables named
in it.
\begin{enumerate}[label=(C\arabic*),leftmargin=*]
\item An \emph{\CondI-word}: a word $w$ with
  $(q,w)\in\mu_\psi$ and $(q,w)\notin\mu_\varphi$ for some $q\in X^\pm$.
\item A \emph{\CondU-word}: a word $w$ such that $(p,w)$ is
  live in $\varphi$ and $\DeadU_\psi(p,w)$ holds (\cref{def:sm-deadu}), for
  some $p\in X^\pm$.
\item A \emph{hole word}, used only when \CondU\ holds: the
  separator word $u\sep t$ of a hole $(u,t)$ of an extracted instance
  $I_{pq}$ with $p,q\in X^\pm$ (\cref{def:ex-instance}).
\end{enumerate}
Its \emph{length} is the length of the word, so a hole word has length
$\len u+\len t+1$.
\end{definition}

By \cref{prop:sm-decide}(2), an \CondI-word exists iff \CondI\ fails. By
\cref{lem:un-dead}, a \CondU-word exists iff \CondU\ fails. If \CondU\ holds,
then by \cref{thm:ex-reduction}(1) a hole word exists iff \CondII\ fails.

\begin{theorem}[Main theorem]\label{thm:as-main}
Let $n\ge0$ be an arity and $c$ a variance of $f$.
\begin{enumerate}[leftmargin=*]
\item Existential entailment $\varphi\models_{\mathcal D}\Exists\psi$ over
  $\BFT$ is decidable for each of the tree classes $\mathcal D=\Trees$,
  $\Trees_{\mathrm{reg}}$ and $\Trees_{\mathrm{fin}}$.
\item If $f$ is covariant, existential entailment over $\BF$ is decidable for
  each of the three tree classes.
\item Let $\varphi$ over $X$ and $\psi$ over $X\cup Z$ be flat systems over
  $\BFT$, with $v=\len X$ and $s=\len X+\len Z$, and let $\mathcal D$ be a tree
  class. If $\varphi$ and $\psi$ are $\mathcal D$-satisfiable and
  $\varphi\not\models_{\mathcal D}\Exists\psi$, then $(\varphi,\psi)$ has a
  counterexample word of length at most $E_3(10v+12s+15)\le E_3(22s+15)$.
  Every counterexample word determines at most two visible observations such
  that $\varphi$ has solutions with values in $\mathcal D$ in which they hold,
  and no such solution has an extension with values in $\mathcal D$ to a
  solution of $\psi$.
\item Let $f$ be covariant, let $\varphi$ and $\psi$ be flat systems over
  $\BF$ with $v=\len X$ and $s=\len X+\len Z$, and let $\mathcal D$ be a tree
  class. If
  $\varphi$ and $\psi$ are $\mathcal D$-satisfiable over $\BF$ and the
  entailment fails over $\BF$ in $\mathcal D$, then the pair
  $(\varphi^\natural,\psi^\natural)$ of \cref{thm:tc-gadget} has a
  counterexample word of length at most $E_3(20v+24s+15)\le E_3(44s+15)$, and
  every counterexample word of this pair determines a solution of $\varphi$
  over $\BF$ in $\mathcal D$ without an extension over $\BF$ in $\mathcal D$ to
  a solution of $\psi$.
\end{enumerate}
The bounds in (3) and (4) depend neither on $n$ nor on $c$.
\end{theorem}

Unflattened inputs are first flattened (\cref{lem:bg-flattening}), and $s$
then counts the variables after flattening, as the bounds require
(\cref{sec:bg-constraints}). The satisfiability hypotheses in (3) and (4) are
needed only to make the certificate a word: if $\varphi$ is not
$\mathcal D$-satisfiable the entailment holds, and if $\varphi$ is
$\mathcal D$-satisfiable and $\psi$ is not, every $\mathcal D$-solution of
$\varphi$ is a countermodel. The bounds concern certificates, not
countermodels or witnesses: a countermodel is in general an infinite regular
tree assignment (\cref{lem:as-countermodel}), and when the entailment holds the witnesses for the hidden variables may have
to be infinite (\cref{rem:tc-examples}). Part (2) is not claimed for $\BF$
with a contravariant argument; the reduction of \cref{thm:tc-gadget} needs
covariance (\cref{rem:tc-gadget-needed}).

\subsection{The Decision Procedure}\label{sec:as-procedure}

The input consists of the arity $n$, the variance $c$, the signature ($\BFT$,
or $\BF$ with $f$ covariant), a tree class $\mathcal D$, and two finite
conjunctions $\varphi$ and $\psi$ of inequalities and equalities between
type expressions, with visible variables $X$ and hidden variables $Z$. The
output is whether $\varphi\models_{\mathcal D}\Exists\psi$ holds and, if it
fails, a certificate: a counterexample word when both systems are
$\mathcal D$-satisfiable, and otherwise the failure of the
$\mathcal D$-satisfiability test for $\psi$. Each step cites the result that
justifies it; \cref{tab:as-steps} lists the sections involved.

\begin{enumerate}[label=\textbf{Step \arabic*.},ref=\arabic*,leftmargin=*,align=left]
\item\label[asstep]{as:step-flatten} \emph{Flatten.} Replace $\varphi$ and
  $\psi$ by the flat systems of \cref{lem:bg-flattening}. The answer is
  unchanged in every tree class and for both signatures. From now on $s$
  counts the variables after flattening.
\item\label[asstep]{as:step-nullary} \emph{Arity zero.} If $n=0$, the trees
  are $\Bot<f()<\Top$, all finite, and the three classes coincide
  (\cref{rem:bg-nullary,rem:tc-nullary}); over $\BF$ only $\Bot$ and $f()$
  remain. Enumerate the at most $3^s$ assignments of $X\cup Z$, check
  $\varphi$ and $\psi$ for each, and answer by \cref{def:bg-existential}.
  Stop.
\item\label[asstep]{as:step-gadget} \emph{The signature $\BF$.} If the
  signature is $\BF$, replace $(\varphi,\psi)$ by
  $(\varphi^\natural,\psi^\natural)$, with visible variables $X\cup Y_X$ and
  hidden variables $Z\cup Y_Z$, and continue over $\BFT$. The answer is
  unchanged in every tree class (\cref{thm:tc-gadget}); the size parameter
  doubles (\cref{cor:tc-bf}).
\item\label[asstep]{as:step-sat} \emph{Satisfiability in $\mathcal D$.}
  Decide whether $\varphi$ and $\psi$ are $\mathcal D$-satisfiable: for
  $\mathcal D=\Trees$ by \cref{prop:sm-decide}(1), for
  $\mathcal D=\Trees_{\mathrm{reg}}$ by the same test, since
  $\Sat_{\mathrm{reg}}=\Sat$ (\cref{thm:bg-nsse}(1)), and for
  $\mathcal D=\Trees_{\mathrm{fin}}$ by \cref{thm:bg-nsse}(2). If $\varphi$ is
  not $\mathcal D$-satisfiable, answer \emph{holds}. If $\varphi$ is
  $\mathcal D$-satisfiable and $\psi$ is not, answer \emph{fails}: every
  solution of $\varphi$ in $\mathcal D$ is a countermodel. Otherwise both
  systems are satisfiable over arbitrary trees; continue.
\item\label[asstep]{as:step-i0} \emph{Condition \CondI.} For each $q\in X^\pm$,
  decide whether some word $w$ has $(q,w)\in\mu_\psi$ and
  $(q,w)\notin\mu_\varphi$ (\cref{prop:sm-decide}(2)). If so, answer
  \emph{fails}, with the \CondI-word $w$.
\item\label[asstep]{as:step-u} \emph{Condition \CondU.} For each
  $p\in X^\pm$, test the words $w$ of length at most
  $\holeBound(N_{\mathrm U},d_{N_{\mathrm U}})$, with
  $N_{\mathrm U}=2^{4\len X+2\len Z+1}$ as in \cref{thm:un-decidable},
  for $(p,w)$ live in $\varphi$ and $\DeadU_\psi(p,w)$, using the source
  automaton $\Bsrc_p$ and its self guard
  (\cref{thm:un-decidable,lem:un-source,lem:un-holes}). If one passes, answer
  \emph{fails}, with the \CondU-word $w$.
\item\label[asstep]{as:step-ii} \emph{Condition \CondII.} Now \CondI\ and
  \CondU\ hold. For each $p,q\in X^\pm$, compute the extracted instance
  $I_{pq}$ (\cref{def:ex-instance}) and
  \[
   K_{pq}=\max\bigl(\Lnm(N_A,N_B,N_{A'},N_{B'}),\ \Lensharp_{I_{pq}}\bigr),
  \]
  with the state counts of $I_{pq}$, and test the pairs $(u,t)$ with
  $\len u+\len t\le K_{pq}$ for conditions (a)--(e) of \cref{def:ex-ss2self}.
  This search finds a hole iff $I_{pq}$ has one: by \cref{lem:ex-fh},
  $I_{pq}$ has forced target heads, so \cref{cor:nm-short-holes} leaves only
  holes shorter than $\Lnm$ and mortal holes, and \cref{thm:mo-short} shortens
  the latter (\cref{thm:mo-decide}; numerically, \cref{cor:mo-extracted}). If
  some $I_{pq}$ has a hole $(u,t)$, answer
  \emph{fails}, with the hole word $u\sep t$; otherwise answer \emph{holds}
  (\cref{cor:ex-entailment}).
\item\label[asstep]{as:step-classes} \emph{Tree classes.} Return the answer.
  \Cref{as:step-i0,as:step-u,as:step-ii} decide entailment over arbitrary
  trees, and their answer is the answer for $\mathcal D$: $\Ereg=\Eall$
  (\cref{thm:tc-regular}), and $\Efin=\Eall$ when both systems have finite
  solutions (\cref{thm:tc-finite}). A countermodel in $\mathcal D$ is obtained
  from a counterexample word by \cref{lem:as-countermodel}. Over $\BF$, a
  countermodel of $(\varphi^\natural,\psi^\natural)$, including one returned
  in \cref{as:step-sat}, is restricted to the original visible variables
  (\cref{lem:as-countermodel}(4)).
\end{enumerate}

\begin{table}[tbp]
\caption{The steps of the decision procedure, the results that justify them,
and the sections where these are proved. The characterization
(\cref{thm:ch-characterization}), built on the leaf cones of \cref{sec:lf},
underlies \cref{as:step-sat,as:step-i0,as:step-u,as:step-ii}.}
\label{tab:as-steps}
\centering
\small
\begin{tabular}{@{}ll>{\raggedright\arraybackslash}p{0.40\linewidth}l@{}}
\toprule
Step & Task & Results used & Sections\\
\midrule
\cref{as:step-flatten} & flatten & \cref{lem:bg-flattening} & \cref{sec:bg}\\
\cref{as:step-nullary} & arity zero & \cref{rem:bg-nullary,rem:tc-nullary} &
  \cref{sec:bg,sec:tc}\\
\cref{as:step-gadget} & $\BF$ & \cref{lem:tc-gadget,thm:tc-gadget,cor:tc-bf} &
  \cref{sec:tc}\\
\cref{as:step-sat} & satisfiability & \cref{prop:sm-decide}(1),
  \cref{thm:bg-nsse,thm:bg-nsse-sat},
  \cref{thm:ch-characterization}(1) &
  \cref{sec:bg,sec:sm,sec:ch}, \cref{app:bg}\\
\cref{as:step-i0} & \CondI & \cref{lem:sm-mu}, \cref{prop:sm-decide}(2) &
  \cref{sec:sm}\\
\cref{as:step-u} & \CondU & \cref{lem:un-holes,lem:un-rejected-path,thm:un-decidable};
  \cref{cor:bg-hole-decision} & \cref{sec:un,sec:bg}\\
\cref{as:step-ii} & \CondII & \cref{thm:ex-reduction,cor:ex-entailment,lem:ex-fh};
  \cref{cor:nm-short-holes}; \cref{thm:mo-short,thm:mo-decide,cor:mo-extracted} &
  \cref{sec:ex,sec:nm,sec:mo}\\
\cref{as:step-classes} & tree classes & \cref{thm:tc-regular,thm:tc-finite,cor:tc-decidable} & \cref{sec:tc}\\
\bottomrule
\end{tabular}
\end{table}

Condition \CondU\ is decided in \cref{as:step-u} before the holes of the
extracted instances are used in \cref{as:step-ii}. This order matters:
\cref{thm:ex-reduction} assumes \CondU, and without it a failure of \CondII\
need not be a hole (\cref{ex:ex-u-needed}). The bounded search of
\cref{as:step-ii} itself needs only forced target heads, condition (FH), which
every extracted instance satisfies by construction (\cref{lem:ex-fh}); without
(FH) the search would be incomplete (\cref{ex:nm-fh}). Equivalently, after \cref{as:step-sat} one may test all
candidate words of the three kinds of \cref{def:as-certificate} up to the
bound of \cref{thm:as-main}(3): a candidate hole of some $I_{pq}$ found while
\CondU\ fails still gives the correct answer \emph{fails}, and if no candidate
passes, all three conditions hold by \cref{prop:as-bounds}.

\subsection{Countermodels}\label{sec:as-countermodels}

A counterexample word determines a countermodel through least closures: the
least model of $\varphi$ that contains the seed of the word is a solution of
$\varphi$ that no solution of $\psi$ extends. The word also gives at most two
observations, which ground bounds pin in every tree class.

\begin{lemma}[Countermodels]\label{lem:as-countermodel}
Let $\varphi$ over $X$ and $\psi$ over $X\cup Z$ be flat systems over $\BFT$,
both satisfiable, and let a counterexample word of $(\varphi,\psi)$ be given.
Define a set $A$ of seeds and a set $\mathcal O$ of visible observations
(\cref{def:tc-observation}) by
\begin{itemize}[leftmargin=*]
\item for an \CondI-word $w$ with $q$: $A=\varnothing$, and $\mathcal O$ denies
  $(q,w)$;
\item for a \CondU-word $w$ with $p$: $A=\{(p,w)\}$, and $\mathcal O$ asserts
  $(p,w)$;
\item for a hole word $u\sep t$ of $I_{pq}$: $A=\{(p,t)\}$, and $\mathcal O$
  asserts $(p,t)$ and denies $(q,u)$.
\end{itemize}
\begin{enumerate}[leftmargin=*]
\item $H_\varphi(A)$ is a model of $\varphi$, namely $\mu_\varphi$,
  $\Cl_\varphi(p,w)$ or $\Cl_\varphi(p,t)$. It decodes to a solution
  $\rho_A$ of $\varphi$ in which the observations of $\mathcal O$ hold.
\item No assignment of $X$ in which the observations of $\mathcal O$ hold
  extends to a solution of $\psi$. In particular, $\rho_A$ is a countermodel
  over arbitrary trees.
\item Let $\mathcal D$ be a tree class in which $\varphi$ is satisfiable, and
  let $B$ be the set of finite ground bounds equivalent to the observations of
  $\mathcal O$ by \cref{lem:tc-pins}. Then $\varphi\cup B$ has a solution in
  $\mathcal D$, and every such solution is a solution of $\varphi$ in
  $\mathcal D$ without any extension to a solution of $\psi$.
\item Suppose that $(\varphi,\psi)=(\varphi_0^\natural,\psi_0^\natural)$ for
  systems $\varphi_0$ over $X_0$ and $\psi_0$ over $X_0\cup Z_0$ over $\BF$,
  with $f$ covariant, as in \cref{thm:tc-gadget}. If $\rho$ is a solution of
  $\varphi$ with values in a tree class $\mathcal D$ that has no extension
  with values in $\mathcal D$ to a solution of $\psi$, then $\rho|_{X_0}$ is a
  solution of $\varphi_0$ over $\BF$ in $\mathcal D$ without an extension over
  $\BF$ in $\mathcal D$ to a solution of $\psi_0$.
\end{enumerate}
\end{lemma}

\begin{proof}
(1) Since $\varphi$ is satisfiable, $\mu_\varphi$ is its least model
(\cref{lem:lf-closure}(2)). A \CondU-word has $(p,w)$ live in $\varphi$ by
definition, and a hole $(u,t)$ of $I_{pq}$ has $(p,t)$ live in $\varphi$ by
\cref{thm:ex-reduction}, which applies because \CondU\ holds; so
$\Cl_\varphi(p,w)$, respectively $\Cl_\varphi(p,t)$, is the least model of
$\varphi$ that contains the seed (\cref{lem:lf-closure}(2)). The observations
hold in $H_\varphi(A)$: $(q,w)\notin\mu_\varphi$ by definition, the seed lies
in its own closure, and $(q,u)\notin\Cl_\varphi(p,t)$ by
\cref{thm:ex-reduction}. By \cref{lem:lf-clauses}, $H_\varphi(A)$ is
$S_{\rho_A}$ for a unique solution $\rho_A$ of $\varphi$.

(2) Let $\rho'$ be an assignment of $X$ in which the observations hold, and
suppose that it extends to a solution of $\psi$. By \cref{lem:lf-clauses},
$S_{\rho'}=T|_X$ for a model $T$ of $\psi$. Since $T$ contains $\Fact_\psi$
and is closed under the rules of $\psi$, it contains $\mu_\psi$, and it
contains $\Reach_\psi(a)$ for each of its atoms $a$ (\cref{lem:lf-closure}(1)).
For an \CondI-word, $(q,w)\in\mu_\psi\subseteq T$ is visible, so
$(q,w)\in S_{\rho'}$, against the denial. For a \CondU-word, $(p,w)\in T$, so
$T\supseteq\Reach_\psi(p,w)$ contains a forbidden root of $\psi$, against the
consistency of $T$. For a hole word, $(p,t)\in T$ and
$(p,t)\to^*_\psi(q,u)$ by \cref{thm:ex-reduction}, so $(q,u)\in T|_X=S_{\rho'}$,
against the denial. Applied to $\rho_A$, this proves the second claim;
\cref{thm:ex-reduction}(2) is the case of a hole word.

(3) By \cref{lem:tc-pins}, an observation holds in an assignment iff the
corresponding ground bound does, so $\rho_A$ solves $\varphi\cup B$. For
$\mathcal D=\Trees$ take $\rho_A$; for $\mathcal D=\Trees_{\mathrm{reg}}$,
\cref{lem:tc-pinning-regular} gives a regular solution, the bounds being
finite and hence regular; and for $\mathcal D=\Trees_{\mathrm{fin}}$,
\cref{thm:tc-pinning} gives a finite one, since $\varphi$ has a finite
solution. Every solution of $\varphi\cup B$ satisfies the observations, so it
has no extension by (2).

(4) The values $\rho(x)$, $x\in X_0$, lie below the values $\rho(y_x)$, which
satisfy $y_x\le f(y_x^n)$, so they are $\Top$-free by \cref{lem:tc-gadget}(1),
and $\rho|_{X_0}$ solves $\varphi_0$ over $\BF$. Suppose that $\sigma$, with
values in $\mathcal D$, extends $\rho|_{X_0}$ to a solution of $\psi_0$ over
$\BF$. Put $\sigma(y_z)=T_{h(\sigma(z))+1}$ for $z\in Z_0$ if
$\mathcal D=\Trees_{\mathrm{fin}}$, and $\sigma(y_z)=\Omega$ otherwise. By
\cref{lem:tc-gadget}(2), $\rho$ together with $\sigma$ solves $\psi$, which
does not constrain the variables $y_x$; this contradicts the choice of $\rho$.
\end{proof}

Informally, a countermodel in $\mathcal D$ can then be found by enumerating
the finite assignments, or finite graph presentations, of $X$ and testing
$\varphi\cup B$; we do not analyze this search.

\begin{example}[Shift closure: the procedure end to end]\label{ex:as-running}
Run the procedure on \cref{ex:intro-shift} with $\varphi_1=\{x\le y\}$ and
$\psi=\{x\le z,\ z\le y,\ z\le f(z,z)\}$, over $\BF$ with $f$ binary and
covariant, in the class of finite trees. Both systems are flat.
\Cref{as:step-gadget} adds $\mathsf{TF}(x,y_x)$ and $\mathsf{TF}(y,y_y)$ to
$\varphi_1$ and $\mathsf{TF}(z,y_z)$ to $\psi$; the assignment $\Bot$ to every
variable solves both new systems, so \cref{as:step-sat} continues.
\begin{itemize}[leftmargin=*]
\item \Cref{as:step-i0}: $\psi^\natural$ has no facts, so
  $\mu_{\psi^\natural}=\varnothing$ and \CondI\ holds.
\item \Cref{as:step-u}: every atom $(v^0,w)$ with $v\in\{x,y,y_x,y_y\}$ is
  dead by support in $\varphi_1^\natural$, since for $v\in\{x,y\}$ the copy
  $v^0\to y_v^0$, and then pops of $y_v^0$, lead to the forbidden root
  $(y_v^0,\eps)$. Since $f$ is covariant, pure moves preserve signs
  (\cref{lem:lf-moves}(2)), and the forbidden roots $(z^0,\eps)$ and
  $(y_z^0,\eps)$ of $\psi^\natural$ have sign $0$, so no atom of sign $1$ is
  root-dead in $\psi^\natural$. Hence no $\RU(p)$ rejects a word, and \CondU\
  holds, as in \cref{ex:un-running}.
\item \Cref{as:step-ii}: the computation of \cref{ex:ex-running} carries over,
  with $y_y$ in the role of $g$. The spine automaton of $x^1$ moves from
  $\{x^1,y^1,y_x^1,y_y^1\}$ to $\{y_x^1,y_y^1\}$ on every letter and is safe;
  $J(t)=0$ for every target $t$; the only possible event is $(0,0)$, present
  iff $t$ is a prefix of $u$; and the head events are the pairs $(s,0)$ with
  $\wsuf us=t$. So $(01,1)$ is a hole of $I_{y^1x^1}$, and the search may
  answer \emph{fails} with the hole word $01\sep1$, of length~$4$.
\end{itemize}
By \cref{lem:as-countermodel}, the word gives $A=\{(y^1,1)\}$, and
$\mathcal O$ asserts $(y^1,1)$ and denies $(x^1,01)$. The least model
$\Cl_{\varphi_1^\natural}(y^1,1)$ is the regular countermodel
$x=y=f(\Omega,\Bot)$, $y_x=y_y=\Omega$: the closure of the seed consists of
the atoms $(y^1,1\sigma)$ and $(x^1,1\sigma)$. For finite trees, the pins of
\cref{lem:tc-pins} for the two observations are
\[
 y\le\emb_\Top(1,\Bot)=f(\Top,\Bot),\qquad
 f(f(\Bot,f(\Bot,\Bot)),\Bot)=\emb_\Bot(01,f(\Bot,\Bot))\le x .
\]
Every finite solution of $\varphi_1^\natural$ with these two bounds is a
countermodel, for instance $x=y=f(f(\Bot,f(\Bot,\Bot)),\Bot)$ with
$y_x=y_y=T_3$, and its restriction to $\{x,y\}$
(\cref{lem:as-countermodel}(4)) is the countermodel of \cref{ex:intro-shift}.
On trees without $\Top$, the two pins say $1\notin\Kpos(y)$ and
$01\in\Kpos(x)$, which is exactly the failure of
$\Suf(\Kpos(x))\subseteq\Kpos(y)$ found there.
\end{example}

\subsection{Counterexample Bounds}\label{sec:as-bounds}

The automata of the extracted instances have
\[
 N_A=2^{2v},\qquad N_B=2^{4v+2h+1}=2^{2v+2s+1},\qquad N_{A'}=N_{B'}=2^{2s}
\]
states (\cref{sec:ex-instances}), and the unit readers of condition \CondU\
have $N_{\mathrm U}=N_B$ states (\cref{thm:un-decidable}).

\begin{proposition}[Counterexample bounds]\label{prop:as-bounds}
Let $\varphi$ over $X$ and $\psi$ over $X\cup Z$ be flat systems over $\BFT$,
both satisfiable, with $v=\len X$ and $s=\len X+\len Z$.
\begin{enumerate}[leftmargin=*]
\item If \CondI\ fails, there is an \CondI-word of length less than
  $2^{2v+2s+2}$.
\item If \CondU\ fails, there is a \CondU-word of length at most
  $\holeBound(N_{\mathrm U},d_{N_{\mathrm U}})\le E_3(10v+10s+17)$.
\item If \CondI\ and \CondU\ hold and \CondII\ fails, there are
  $p,q\in X^\pm$ and a hole $(u,t)$ of $I_{pq}$ such that either
  $\len u+\len t\le\Lnm(N_A,N_B,N_{A'},N_{B'})$, a bound singly exponential in
  $s$, or $(u,t)$ is a mortal hole and
  $\len{u\sep t}\le\Lensharp_{I_{pq}}\le E_3(10v+12s+15)$.
\end{enumerate}
If $v\ge1$, each of these bounds is at most $E_3(10v+12s+15)\le E_3(22s+15)$,
and the mortal closed form is the largest. If $v=0$, there is no counterexample
word.
\end{proposition}

\begin{proof}
(1) This is \cref{prop:sm-decide}(2), with $\len X=v$ and $\len{X\cup Z}=s$.

(2) This is \cref{thm:un-decidable}.

(3) By \cref{thm:ex-reduction}(1), which applies because $\varphi$ is
satisfiable and \CondU\ holds, some $I_{pq}$ has a hole. It satisfies (FH) by
\cref{lem:ex-fh}, so by \cref{cor:nm-short-holes} it has a hole with
$\len u+\len t\le\Lnm(N_A,N_B,N_{A'},N_{B'})$ or a mortal hole. The function
$\Lnm$ is a polynomial in the four state counts (\cref{sec:nm-three-cases}),
each of which is at most $2^{2v+2s+1}$, so the first bound is singly
exponential in $s$. In the second case, \cref{thm:mo-short} gives a mortal hole
whose separator word has length at most $\Lensharp_{I_{pq}}$, and
\cref{cor:mo-extracted} bounds this by $E_3(10v+12s+15)$.

For the comparison, let $v\ge1$, so that $s\ge1$. The bound of (1) is at most
$2^{4s+2}$, and the non-mortal hole word has length at most
$\Lnm+1\le54N_A^2N_BN_{A'}N_{B'}+1<2^{6v+6s+7}$ (\cref{sec:nm-three-cases});
both are far below the triply exponential bounds. Of these,
$E_3(10v+10s+17)\le E_3(10v+12s+15)$ since $s\ge1$, so the mortal closed form is the
largest, and $E_3(10v+12s+15)\le E_3(22s+15)$ since $v\le s$. If $v=0$, then
$X^\pm$ is empty, and \cref{def:as-certificate} admits no word.
\end{proof}

\begin{table}[tbp]
\caption{Length bounds for shortest counterexample words
(\cref{def:as-certificate}), with $v=\len X$ and $s=\len X+\len Z$, and their
growth in $s$: singly or triply exponential. Over $\BF$ the reduction of
\cref{thm:tc-gadget} replaces $v$ and $s$ by $2v$ and $2s$.}
\label{tab:as-bounds}
\centering
\small
\begin{tabular}{@{}lllll@{}}
\toprule
Failing condition & Word & Length bound & Growth & Source\\
\midrule
\CondI & $w$ & $<2^{2v+2s+2}$ & single &
  \cref{prop:as-bounds}(1)\\
\CondU & $w$ & $\le E_3(10v+10s+17)$ & triple &
  \cref{thm:un-decidable}\\
\CondII, not mortal & $u\sep t$ & $\le\Lnm(N_A,N_B,N_{A'},N_{B'})+1$ & single &
  \cref{cor:nm-short-holes}\\
\CondII, mortal & $u\sep t$ & $\le E_3(10v+12s+15)$ & triple &
  \cref{thm:mo-short,cor:mo-extracted}\\
\bottomrule
\end{tabular}
\end{table}

\Cref{tab:as-bounds} summarizes the cases. Both triply exponential terms come
from the bounded-depth search of \citet{Zhao26} and its anchored
extension, and in both cases the depth is bounded through a short rejected
tail (\cref{thm:bg-depth}). A mortal hole has a rejected tail shorter than
$N_A=2^{2v}$, which involves only visible variables; its reader has
$\Nbar\le2^{2v+4s+2}$ states, and its depth is at most
$d_I=2\Nbar N_A-1<2^{4v+4s+3}$ (\cref{lem:mo-depth}). A unit reader has
$N_{\mathrm U}=2^{2v+2s+1}$ states, and the search of
\cref{cor:bg-hole-decision} bounds the depth of a suitable hole by
$d_{N_{\mathrm U}}<2^{4v+4s+5}$. The mortal closed form is the largest of the
closed forms in the table. This does not order the exact bounds: for
$s\le2$ and $v\ge1$, the exact \CondU\ bound exceeds the exact mortal bound,
because $d_{N_{\mathrm U}}\approx8N_{\mathrm U}^2$ is much larger than $d_I$.
In the part that is not mechanized, the smaller reader of
\cref{sec:un-smaller} gives an even smaller bound for \CondU\
(\cref{prop:un-sharper}), and \cref{sec:cx} analyzes the complexity of the
procedure.

\subsection{Proof of the Main Theorem}\label{sec:as-proof}

\begin{proof}[Proof of \cref{thm:as-main}]
\emph{Parts (1) and (2).} Each step of the procedure of
\cref{sec:as-procedure} terminates and is correct by the result it cites
(\cref{tab:as-steps}); three points need a remark. First, if the procedure
passes \cref{as:step-sat}, both systems are satisfiable in $\mathcal D$, hence
over arbitrary trees, so \cref{as:step-i0,as:step-u,as:step-ii} decide
entailment over arbitrary trees by \cref{thm:ch-characterization}(2), and
\cref{as:step-classes} transfers the answer to $\mathcal D$. Second,
\cref{as:step-u} runs before \cref{as:step-ii}, so \CondU\ holds when
\cref{as:step-ii} starts, and \cref{thm:ex-reduction,thm:mo-decide} apply to
every $I_{pq}$, which has forced target heads by \cref{lem:ex-fh}. Third,
\cref{as:step-nullary,as:step-u,as:step-ii} are finite searches: there are
finitely many assignments at arity zero, the alphabet $\Sigma$ is finite, the
bounds are computable, and each candidate is tested over finitely many cuts.
This proves (1), and (2) follows by \cref{as:step-gadget}, which preserves the
answer by \cref{thm:tc-gadget}.

\emph{Part (3).} Let $\varphi$ and $\psi$ be $\mathcal D$-satisfiable with
$\varphi\not\models_{\mathcal D}\Exists\psi$. Then both are satisfiable over
arbitrary trees, and $\varphi\not\models\Exists\psi$: for
$\mathcal D=\Trees_{\mathrm{reg}}$ by \cref{thm:tc-regular}, and for
$\mathcal D=\Trees_{\mathrm{fin}}$ by \cref{eq:tc-finite}, since both systems
have finite solutions. By \cref{thm:ch-characterization}(2), \CondI, \CondU\
or \CondII\ fails. In the first two cases \cref{prop:as-bounds}(1,2) gives an
\CondI-word or a \CondU-word; otherwise \CondI\ and \CondU\ hold and
\cref{prop:as-bounds}(3) gives a hole word. A counterexample word exists only
if $v\ge1$, and then its length is at most
$E_3(10v+12s+15)\le E_3(22s+15)$ by \cref{prop:as-bounds}. This argument uses
no hypothesis on $n$; for $n=0$ the only word is $\eps$, and every
counterexample word has length at most $1$. By
\cref{lem:as-countermodel}(3), every counterexample word yields a
countermodel in $\mathcal D$, since $\varphi$ is $\mathcal D$-satisfiable.
None of the bounds involves $n$ or $c$.

\emph{Part (4).} Let $f$ be covariant, and let $\varphi$ and $\psi$ over $\BF$
be $\mathcal D$-satisfiable over $\BF$, with failing entailment over $\BF$ in
$\mathcal D$. For a $\Top$-free tree $t$ in $\mathcal D$, let $y(t)$ be the
tree of \cref{lem:tc-gadget}(2) in $\mathcal D$ with $t\le y(t)\le
f(y(t)^n)$. Extending a $\BF$-solution of $\varphi$ by $y_x=y(\rho(x))$, and a
$\BF$-solution $\sigma$ of $\psi$ by $y_z=y(\sigma(z))$, shows that
$\varphi^\natural$ and $\psi^\natural$ are $\mathcal D$-satisfiable over $\BFT$,
and by \cref{thm:tc-gadget} the entailment of $(\varphi^\natural,
\psi^\natural)$ fails in $\mathcal D$. This pair has $2\len X$ visible and
$2\len Z$ hidden variables (\cref{cor:tc-bf}), so (3), applied with $2v$
and $2s$, gives a counterexample word of length at most
$E_3(10\cdot2v+12\cdot2s+15)=E_3(20v+24s+15)\le E_3(44s+15)$. It yields a countermodel for
$(\varphi^\natural,\psi^\natural)$ in $\mathcal D$, which restricts to a
countermodel over $\BF$ in $\mathcal D$ by \cref{lem:as-countermodel}(4).
\end{proof}

\subsection{Plain Entailment}\label{sec:as-plain}

Plain entailment is the special case of existential entailment without hidden
variables and with the consequent $x\le y$ (\cref{def:bg-existential}). Its
decidability, proved by \citet{Zhao26}, is therefore a corollary, with the same
procedure.

\begin{corollary}[Plain entailment]\label{cor:as-plain}
For every arity $n\ge0$ and every variance $c$, plain entailment
$\varphi\models_{\mathcal D}x\le y$ over $\BFT$ is decidable for each tree class
$\mathcal D$, and
\[
 \varphi\models_{\mathrm{reg}}x\le y\iff\varphi\models x\le y,\qquad
 \varphi\models_{\mathrm{fin}}x\le y\iff
 \neg\Sat_{\mathrm{fin}}(\varphi)\ \lor\ \varphi\models x\le y .
\]
If $f$ is covariant, plain entailment over $\BF$ is decidable for each tree
class.
\end{corollary}

\begin{proof}
With $Z=\varnothing$ and $\psi=\{x\le y\}$, \cref{thm:as-main}(1,2) decides the
entailment. The consequent has the solution $x=y=\Bot$ in every tree class, so
\cref{thm:tc-regular} gives the first equivalence and \cref{eq:tc-finite} the
second.
\end{proof}
